\chapter{模型与表示：从基线到 PyTorch、Transformer 与文本适配}
\label{ch:lower-ml-model}

\begin{itemize}
  \item 直接先修：上册第 5、6、9、10、12、14、15、16、17 章。
  \item 学习产出：能从 NumPy 基线迁移到可复现的 PyTorch 训练，并证明序列模型确实保留顺序。
  \item 研究边界：本章的小模型用于暴露训练与信息协议，不声称达到生产模型容量。
\end{itemize}

\section{先冻结简单基线}

回归树桩在候选阈值 $c$ 上最小化两侧平方误差。候选集合必须来自排序后的唯一特征值中点；若直接使用输入相邻项，中点集合会依赖行顺序。透明实现先计算
\[
\mathcal C=\left\{\frac{x_{(j)}+x_{(j+1)}}2:x_{(j)}<x_{(j+1)}\right\},
\]
再逐个比较。成熟库对照使用 \texttt{DecisionTreeRegressor(max\_depth=1)}；两者在冻结数据上的预测差为 0。两轮残差 boosting 也与 \texttt{GradientBoostingRegressor} 对拍。这里的“双实现”只验证同一平方损失与树深，不代表不同库默认值天然等价。

线性、树和 boosting 不是被深度模型淘汰的装饰品。它们提供训练时间、解释难度和样本外损失的下界；复杂模型若不能在冻结验证协议下改善决策结果，应保留基线。资产定价中的系统比较也表明，模型灵活性必须与正则化和严格样本外评价一起解释\cite{gu2020mlassetpricing}。

\section{PyTorch：从张量到 checkpoint}

张量同时携带数值、形状、dtype 和 device。训练循环的最小闭环分为三步：
\begin{enumerate}
  \item \texttt{Dataset} 保存样本语义，\texttt{DataLoader} 决定批次与顺序；
  \item 前向传播产生损失，\texttt{autograd} 计算梯度，优化器更新参数；
  \item 在评估模式与无梯度上下文中冻结推理。
\end{enumerate}
代码还固定随机种子、CPU device、epoch 数与学习率，并把配置和参数字节 $D$ 哈希为 checkpoint 指纹 $h(D)$\cite{paszke2019pytorch}。

对四个样本训练两层 MLP，最终 MSE 约为 $6.34\times10^{-4}$。这个数不是模型质量结论；真正的证据是同一配置两次运行得到完全相同预测和指纹。若更换 GPU、算子或框架版本，必须重新声明可复现级别，而非默认逐位一致。

\begin{diagnostic}
只保存 \texttt{state\_dict} 不够：还要保存 tokenizer、特征顺序、预处理参数、随机种子、优化器状态、代码版本和数据截止。checkpoint 是推理契约，不只是权重文件。
\end{diagnostic}

\section{因果卷积、RNN、Attention 与 Transformer}

序列输入为 $X\in\mathbb R^{B\times T\times F}$。因果卷积的 $t$ 时点只读取 $\{X_s:s\le t\}$；RNN 递推 $h_t=\phi(W_xX_t+W_hh_{t-1})$；self-attention 通过
\[
\operatorname{softmax}\!\left(\frac{QK^{\mathsf T}}{\sqrt d}+M\right)V
\]
聚合位置，其中因果 mask $M$ 屏蔽未来。Transformer 再叠加前馈、残差与归一化\cite{vaswani2017attention}。没有位置编码的 attention 对置换近似等变，不能自动识别先后。

冻结负例把 $[1,2,4,8]$ 反转。时间均值差为 0，而因果卷积、RNN、带位置与因果 mask 的 Attention/Transformer 输出均发生变化。这个实验不证明模型预测有效，只证明实现没有把顺序悄悄压成均值。padding mask、缺失 mask 与因果 mask 必须分别记录。

\section{NLP/LLM：任务边界先于模型规模}

文本链从 tokenizer 和 embedding 开始，再进入预训练编码器、任务头与决策映射。零样本依赖任务提示或词义规则；少样本利用少量标注形成原型；全量微调更新编码器全部参数；PEFT 只训练 adapter、LoRA 或任务头。LoRA 的核心是冻结预训练权重并学习低秩更新\cite{hu2022lora}。冻结实验以小型编码器代理说明参数边界：可训练参数仅占总参数约 $15.8\%$，不是把它包装成基础模型。

真实预训练编码器必须绑定模型卡、许可、权重版本和 tokenizer。新闻或公告还要同时保存首次发布时间、抓取时间与修订时间；决策日后的修订不能回填历史。LLM 生成的标签需保留 prompt、采样参数、原文证据和人工抽查，零/少样本结果不能免除时间审计。

\begin{implementationnote}
\begin{lstlisting}[language=bash]
uv run python notebooks/lower/ml_alpha_model.py \
  evidence/ml-alpha-model/oracle.json
\end{lstlisting}
\end{implementationnote}

本单元向 Route Capstone 交付：基线表、PyTorch checkpoint、顺序负例、文本适配协议和算力/随机性限制。

\section{分层练习}
\begin{enumerate}[leftmargin=*]
  \item \textbf{口述概念：}tensor 的 shape、dtype、device 和梯度状态分别约束什么？
  \item \textbf{口述概念：}为什么树桩候选阈值必须先排序唯一特征值？
  \item \textbf{笔试推导：}写出两层 MLP 的前向传播与平方损失梯度链。
  \item \textbf{笔试推导：}解释没有位置编码的 self-attention 对时间置换为何不敏感。
  \item \textbf{数值编程：}把序列反转，比较均值、因果卷积、RNN、Attention 与 Transformer 输出。
  \item \textbf{数值编程：}保存并恢复 checkpoint，验证预测与哈希不漂移。
  \item \textbf{研究判断：}审计一个把随机特征均值池化称作“深度序列模型”的错误研究包。
  \item \textbf{研究判断：}比较全量微调、只训任务头和 LoRA/adapter 的证据边界。
\end{enumerate}

